\documentclass[../main.tex]{subfiles}

\begin{document}
\subsection{Algoritma, {\Pseudocode} dan {\Flowchart}}
Algoritma adalah spesifikasi urutan langkah untuk melakukan pekerjaan tertentu. Dalam bahasa Indonesia, definisi algoritma adalah logika, metode dan tahapan (urutan) sistematis yang digunakan untuk memecahkan suatu permasalahan \parencite{utami10}.

{\Pseudocode} adalah bentuk informal untuk mendeskripsikan algoritma yang mengikuti struktur bahasa pemrograman tertentu dengan tujuan untuk memudahkan dibaca oleh manusia, mudah untuk dipahami dan mudah dalam menuangkan ide atau hasil pemikiran \parencite{harumy2016belajar}.

{\Flowchart} didefinisikan sebagai skema penggambaran dari algoritma atau proses \parencite{harumy2016belajar}. Tabel berikut menampilkan simbol-simbol yang digunakan dalam menyusun {\flowchart} \parencite{samala2021pemrograman}:
{\cpsize Halo}

\inputtable{|C|l|}{
    Simbol & {\hspace{30ex}Keterangan} \\ \hline
    {\vspace{0.2ex}
\includegraphics[scale=0.398]{res/m1i2.png}} & Simbol permulaan awal ({\sstart}) dan akhir program ({\sstop}). \\ \hline
    {\vspace{0.2ex}\includegraphics[scale=0.398]{res/m1i3.png}} & Simbol yang digunakan untuk menghubungkan antar simbol. \\ \hline
    {\vspace{0.2ex}\includegraphics[scale=0.398]{res/m1i4.png}} & Simbol yang menunjukkan proses dalam {\flowchart}. \\ \hline
    {\vspace{0.2ex}
\includegraphics[scale=0.398]{res/m1i5.png}} & Simbol pemilihan proses, pengambilan keputusan berdasarkan suatu kondisi tertentu. \\ \hline
    {\vspace{0.2ex}\includegraphics[scale=0.398]{res/m1i6.png}} & Simbol yang menyatakan proses {\ioi} dan {\ioo}. \\ \hline
    {\vspace{0.2ex}\includegraphics[scale=0.398]{res/m1i7.png}} & Simbol untuk pemanggilan suatu prosedur. \\ \hline
    {\vspace{0.2ex}
\includegraphics[scale=0.398]{res/m1i8.png}} & Simbol penghubung dalam satu halaman. \\ \hline
    {\vspace{0.2ex}
\includegraphics[scale=0.398]{res/m1i9.png}} & Simbol penghubung dalam halaman yang berbeda. \\ \hline
    {\vspace{0.2ex}\includegraphics[scale=0.398]{res/m1i10.png}} & Persiapan atau pemberian nilai awal atau proses inialisasi. \\ \hline
    {\vspace{0.2ex}
\includegraphics[scale=0.398]{res/m1i11.png}} & Simbol yang menunjukkan pengolahan operasi secara manual. \\ \hline
    {\vspace{0.2ex}\includegraphics[scale=0.398]{res/m1i12.png}} & Simbol yang menunjukkan {\ioi} manual langsung melalui {\keyboard}. \\ \hline
}{Simbol-Simbol dalam {\Flowchart}}{tab:m1:01}

% \subsection{Bahasa Pemrograman}
Program adalah kumpulan instruksi komputer, program ditulis dengan menggunakan bahasa pemrograman. Jadi bisa disebut bahwa program adalah implementasi dari bahasa pemrograman \parencite{utami10}. Dimana bahasa pemrograman merupakan teknik komando atau instruksi standar untuk memerintah komputer \parencite{dewi2010media}.

\end{document}
